% ECONOMICS_PROBLEM: {"id": "G21-258", "title": "Microcredit heterogeneity", "jel_code": "G21", "source_id": "OP-0373"}
% !TeX program = xelatex
\documentclass[11pt,a4paper]{article}
\usepackage[margin=23mm]{geometry}
\usepackage{fontspec}
\setmainfont{Latin Modern Roman}
\usepackage{amsmath,amssymb,mathtools}
\usepackage{xcolor,enumitem,booktabs,longtable,array,xurl,hyperref}
\definecolor{muted}{HTML}{53636C}
\hypersetup{colorlinks=true,urlcolor=blue,linkcolor=blue}
\setlength{\parindent}{0pt}
\setlength{\parskip}{5pt}
\newcommand{\R}{\mathbb R}
\newcommand{\E}{\mathbb E}
\newcommand{\N}{\mathbb N}
\newcommand{\ind}{\mathbf 1}
\newcommand{\Var}{\operatorname{Var}}
\newcommand{\Cov}{\operatorname{Cov}}
\newcommand{\supp}{\operatorname{supp}}
\DeclareMathOperator*{\argmin}{arg\,min}
\DeclareMathOperator*{\argmax}{arg\,max}
\newcommand{\entrymeta}[3]{{\sffamily\small\color{muted}\textbf{\texttt{#1}}\quad|\quad #2\par #3\par}}
\newcommand{\Problemref}[1]{\texttt{#1}}
\newcommand{\JELref}[2]{\texttt{#2} (JEL \texttt{#1})}
\begin{document}
\section*{Microcredit heterogeneity}
\textbf{Problem ID:} \texttt{G21-258}\quad \textbf{JEL:} \texttt{G21}\par
% BEGIN ECONOMICS STATEMENT
\label{problem:OP-0373}
\label{audited:ECON-085}
{\sffamily\small\color{muted}Original subject group: Development and poverty\par}
\label{definitions:problem:30007655}
\textbf{Audit classification:} Background; proposed extension. Original primary PDF read. NBER landing/PDF access failed; the cited July 2024 revision was not inspected. Existing entrepreneurial heterogeneity findings are preserved; contract and scale extensions are editorial.
\subsection*{Definitions and precise research target}
\textbf{Complete editorial problem.} Fix a population of potential small-business borrowers and a lender with a declared budget. Contract \(p\) specifies credit access, loan size, interest, repayment timing, collateral, and insurance. Let \(D_i(p)\) indicate borrowing, \(\Pi_{iT}(p)\) business profits at horizon \(T\), and \(C_{iT}(p)\) household consumption. Define intention-to-treat effects \(\tau_\Pi(p,x)=E[\Pi_{iT}(p)-\Pi_{iT}(0)\mid X_i=x]\), where \(X_i\) records baseline business activity, assets, and risk. Measure net borrower welfare, lender losses, and default alongside profits. The research task is to identify which observable borrower groups and contract features generate sustained productive gains and which mainly smooth consumption. Offers can be randomized, but borrower-only effects require additional compliance assumptions; conditioning on take-up changes the population. Distinguish learning, indivisible investment, risk reduction, and selection through appropriate variations. Estimate whether scalable selection rules retain gains when repayment incentives and product-market competition respond. Some microcredit effects and heterogeneous responses have already been documented. This is a proposed remaining design and transportability question; the original supporting results are not proof that a precise universal conjecture or first problem statement exists.

\textbf{Definitions and admissible scope.} A loan offer includes enforcement and default consequences as well as the stated contract. Define baseline groups before access is assigned, so a business-existing-at-baseline indicator is admissible but post-offer borrowing is not an ordinary conditioning variable. Profits deduct depreciation, capital cost, and owner labor under a fixed valuation convention; consumption is per person in real units. Borrower welfare uses declared risk preferences and includes debt liabilities. The offer effect is identifiable under random assignment and consistency; a treatment-on-borrowers interpretation needs exclusion, monotonicity, and relevance or another explicit selection model. A targeting rule is evaluated on all eligible borrowers, including default and displacement costs, and does not rank people solely by estimated realized profits. Fix \(T<\infty\), comparator contract \(p=0\), and a probability law for baseline eligible households. Define \(\tau_C(p,x)=E[C_{iT}(p)-C_{iT}(0)\mid X_i=x]\). If welfare is optimized, supply \(V_i(p)=E[\sum_{t=0}^T\delta^t u_i(c_{it},h_{it},a_{it})]\), finite expectations, real resource costs, lender solvency constraints, and weights; profits are not automatically that utility criterion. A randomized observable targeting rule chooses contract probabilities based only on \(X_i\). Its value includes all eligible households, market spillovers, and defaults under the complete regime. The substantive output is the identified set of subgroup gains and feasible rule rankings, conditional on the contract, selection, interference, and measurement restrictions.
\subsection*{Variants and assumption changes}
\begin{enumerate}
\item Entrepreneur type (source-based scope): distinguish households already operating a business before credit arrived from those without one. The cited Hyderabad results answer an important version of this heterogeneity question; further populations require transport assumptions.
\item Contract design (editorial): randomize repayment grace periods, collateral, and insurance with fixed expected subsidy. Compare durable net utility and lender viability and test whether the best contract differs by baseline capital and risk.
\end{enumerate}
\subsection*{References and source audit}
\textbf{Support and access.} Original primary PDF read. NBER landing/PDF access failed; the cited July 2024 revision was not inspected. Existing entrepreneurial heterogeneity findings are preserved; contract and scale extensions are editorial. \textbf{Locator:} MIT-hosted NBER 26346, October 2019 version, abstract.
\begin{enumerate}\raggedright
\item\hypertarget{definitions:source:30007655:1}{} Abhijit Banerjee; Emily Breza; Esther Duflo; Cynthia Kinnan. 2019. \emph{Can Microfinance Unlock a Poverty Trap for Some Entrepreneurs?}. October 2019 NBER version, abstract.
\url{https://economics.mit.edu/sites/default/files/publications/w26346.pdf}.
DOI: \url{https://doi.org/10.3386/w26346}.
\end{enumerate}

\subsection*{Integrated refinement: ECON-085}
{\sffamily\small\color{muted}Loan flexibility for investment versus consumption needs\par}
This refinement adopts the additional source conventions
(page~\pageref{audited:conventions}), subject to its local restrictions.
\textbf{Purpose, repayment flexibility and a viable contract menu.}
For project cash flows $R_{it}$ and consumption needs $N_{it}$, let $a$ specify
repayment timing, grace periods and default. Compare borrower utility $U_i(a)$
and lender net return $\Pi_i(a)$, maximizing induced borrower welfare subject to
declared lender viability. Specify date-by-date budgets, limited liability,
recovery, discounting and outside options. Menus need type-by-type participation
and self-selection; moral hazard needs action incentives and a response selector.
A single contract choice is not an oracle allocation to unobserved types.
Loan purpose may be endogenous because funds are fungible. Distinguish
productive investment, shock insurance and consumption smoothing, and compare
contracts on a common expected fiscal subsidy. Count principal, funding,
processing and default costs once; aggregate and borrower-specific viability
are different restrictions.
\subsection*{Supplied source audit for ECON-085}
Local [S1], [S2], etc. in this audit and the references immediately below
belong to ECON-085, independently of the original entry's reference numbering.
[S1] and [S2] directly support loans tailored to investment versus consumption. The menu/incentive problem is editorial and overlaps contract theory; the review supplies no universally optimal flexible contract.
\subsection*{References for integrated refinement ECON-085}
{\footnotesize\raggedright
\par\noindent\textbf{[S1]} Jing Cai, Muhammad Meki, Simon Quinn, Erica Field, Cynthia Kinnan, Jonathan Morduch, Jonathan de Quidt, and Farah Said (2025). \emph{Microfinance}. VoxDevLit 3(3), January 2025.\par \textit{Verified locator / source role:} Primary publisher page checked for review identity, authors, version and background. The specific published agenda is corroborated in [S2]; the exact formal target is editorial.\par \url{https://voxdev.org/voxdevlit/microfinance}\par\smallskip
\par\noindent\textbf{[S2]} Oliver Hanney / VoxDev (living compilation). \emph{Evidence gaps: Key unanswered questions in development economics}. Accessed 8 October 2026. The page currently covers 20 topics; the ``16-topics URL is a historical slug.\par \textit{Verified locator / source role:} Topic heading: Microfinance. Supports the broad research agenda; this is not a verbatim quotation or an attribution of the displayed mathematical target.\par \url{https://voxdev.org/topic/evidence-gaps-key-unanswered-questions-across-16-topics-development-economics}\par\smallskip
}

\subsection*{Common conventions: Source mathematical conventions}

These conventions apply to each entry unless explicitly replaced.

\textbf{Models and identification.} A primitive model \(m\) specifies all probability laws, feasible choices, information, equilibrium and measurement rules used by its target. The admissible class \(\mathcal M\) is fixed independently of outcomes. If \(P_O(m)\) is its observable probability law and \(\tau(m)\) the target, its identified set at observable law \(P\) is
\[
 \mathcal I_\tau(P)=\{\tau(m):m\in\mathcal M,\ P_O(m)=P\}.
\]
Point identification means that this set is a singleton. An empty set signals incompatibility of data and assumptions. Coordinate bounds need not describe the joint set. Bounds are sharp only if every retained target is attainable or arbitrarily approachable by compatible models. Finite bounds require explicit restrictions on unobserved quantities; unrestricted confounding, hidden wealth, extrapolation or arbitrary equilibrium selection can yield uninformative sets.

\textbf{Causal interventions.} A potential outcome \(Y_i(p)\) is the outcome under a complete policy regime \(p\), including timing, financing, information and market spillovers. The target population and units are fixed before treatment. With interference, use the complete assignment vector or a justified exposure mapping; individual no-interference notation alone cannot identify an economy-wide intervention. A difference of mean potential outcomes does not require the cross-world joint distribution. Conditioning on post-treatment migration, survival, enrollment or employment changes the estimand and needs separate justification.

\textbf{Statistical statements.} An estimator, confidence set, forecast or test specifies a sampling law, model class and loss/coverage criterion. Uniform \((1-\alpha)\)-coverage means \(\inf_{m\in\mathcal M}\Pr_m\{\tau(m)\in C_n\}\geq1-\alpha\), or an explicitly asymptotic version. The whole parameter space trivially has coverage, so informativeness requires a stated width, risk or power objective. A forecasting improvement is relative to a named benchmark, information set and evaluation population.

\textbf{Equilibrium and optimization.} An equilibrium comprises feasible plans, optimality, beliefs where needed, budget constraints and all market/resource clearing. The primitives or a stated benchmark define each correspondence \(\mathcal E(p;m)\). Nonemptiness is assumed or proved, never inferred just from compact policy sets. A selected-equilibrium target uses a declared selection rule; a robust target quantifies over all permitted equilibria. Use suprema or infima unless attainment is established. An \(\varepsilon\)-optimal policy is feasible and within \(\varepsilon\) of the finite optimum value. Report an empty feasible set separately.

\textbf{Welfare and accounting.} Normative utility, interpersonal normalization, social weights, numeraire, discounting and the use of public receipts are primitives. Behavioral utility need not coincide with the welfare criterion. Household welfare includes ownership income and rents once; adding profits again to compensating variation generally double-counts them. A monetary equivalent is defined by an explicit transfer and price convention, with monotonicity and range conditions. Average effects, mechanism contributions, transfers and welfare losses are different targets. Infinite sums require convergence and dynamic feasible plans require terminal or no-Ponzi conditions.

\textbf{Source versions.} A working-paper date, revision date and journal publication year are distinguished. A DOI for a journal version is not automatically the DOI of an earlier manuscript. Locators refer to the stated version and printed pages where specified. A metadata-only check does not verify a claim about a full-text conclusion. Every entry's source subsection narrows the provenance claim accordingly.

\subsection*{Common conventions: Additional-source status and mathematical conventions}

\label{audited:conventions}

Of the 100 supplied ECON entries, 65 distinct problems appear here; 32 close
matches refine earlier entries, and three duplicate questions map to existing
entries. Appendix~\ref{app:audited-crosswalk} lists every destination.
The attachment's P/R/S/U distinctions and source audits are retained. Its
precise mathematical formalizations are editorial unless the individual audit
says otherwise. Candidate subsidy targets ECON-004--006 retain their U flags;
the resolved approval-core question ECON-007 updates OP-0202 above.
The following supplied conventions govern both this part and the attributed
ECON refinements elsewhere in the catalogue, subject to local restrictions.

\textbf{Parameterized questions.} A problem family lists inputs that must be fixed before an instance is evaluated: population, horizon, policies, information, data law, model class and welfare weights. These are required choices, not quantities the citations are assumed to specify. Undefined applications should not be treated as identified estimands.

\textbf{Indices and integrability.} Sets of agents, firms and regions are finite unless a population measure is explicitly used. \(i,j\) index units, \(r\) regions, \(t\) dates and \(h\) horizons. \(\mathbb E\) and \(\Pr\) refer to the declared population and law; \(\mathbf1\{A\}\) is an indicator. Every displayed expectation and discounted sum needs existence in the stated domain. Infinite horizons use a discount in \((0,1)\) and a convergence condition; finite horizons may use discount one.

\textbf{Optimization and existence.} A displayed \(\max\), \(\min\), or \(\arg\max\) is conditional on attainment. For finite feasible menus this is immediate; general spaces need continuity and compactness/coercivity or another existence argument. Without attainment, the target is the supremum/infimum and arbitrarily near-optimal feasible choices. \(\inf\varnothing=+\infty\). Empty feasibility and an unidentified target are different issues.

\textbf{Potential outcomes.} \(Y_i(z)\) is the outcome under a specified intervention, not the observed conditional mean among people choosing \(z\). In binary comparisons \(\Delta Y_i=Y_i(1)-Y_i(0)\). Specify assignment, treatment versions, compliance, information and observation horizon. Interference is allowed under a full policy vector; compressed exposure notation needs a justified mapping. No interference, consistency and overlap are substantive assumptions, not automatic consequences of notation.

\textbf{Populations and observation.} Fix baseline eligibility and population weights. Conditioning on post-policy employment, survival, relocation or program uptake can change the population being compared. Death is not ordinary loss to follow-up; outcomes truncated by death need a composite convention, joint survival target, or explicitly defined survivor stratum. Fertility-changing interventions need a population functional rather than an unexplained fixed descendant cohort. Measurement and missingness models must be supplied.

\textbf{Identification.} A structural model \(m\in\mathcal M\) implies observable law \(P_m\) and target \(\theta(m)\). Identification means
\[
 P_{m_1}=P_{m_2}\ \Longrightarrow\ \theta(m_1)=\theta(m_2)
 \quad\text{for all }m_1,m_2\in\mathcal M .
\]
Without identification, the target is
\[
 \Theta(P;\mathcal M)=\{\theta(m):m\in\mathcal M,\ P_m=P\}.
\]
Writing \(\theta(P)\) before establishing identification can hide incompatible latent models. Confidence sets additionally need a sampling law, confidence level, asymptotic sequence and uniformity class. Point identification, coverage and informative precision are separate properties.

\textbf{Mediation.} Total effects of randomized assignment are distinct from cross-world natural indirect effects. Belief/effort-channel effects require a meaningful mediator intervention and additional measurement, exclusion or mediation assumptions. Randomization of the initial policy alone does not supply those assumptions. Report bounds or interventional alternatives when the stronger target is unsupported.

\textbf{Equilibrium.} A counterfactual \(W(z)\), \(p(z)\), or \(\mu_t^z\) requires individual optimization, feasibility, government/resource budgets, market clearing and state-transition laws. State equilibrium existence and selection, or report ranges over compatible equilibria. Initial-state integration includes subsequent endogenous transitions; current-state integration must use the policy-dependent distribution. Reduced-form invariance after policy is not assumed.

\textbf{Welfare accounts.} Scores, utility, health, dollars and ecological quantities require explicit conversion before addition. Fees, taxes, subsidies, wages and extortion payments are transfers between parties; distributional weights can make them welfare relevant, but their face value is not automatically a resource loss. State ownership, geographic boundaries, foreign-income weights and fiscal finance. Do not add profits to household welfare already including dividends, or count land capitalization and the underlying benefit twice. A benefit-cost expression must distinguish gross benefits from net welfare and subtract every real cost exactly once.

\textbf{Derivatives and risk.} Log derivatives require positive arguments; local responses require admissible perturbations and specified equilibrium branches. Differentiation under expectations needs regularity/domination, and boundaries may allow only one-sided derivatives. Risk uses a specified law; ambiguity uses a declared nonempty law set on a common history space. Adaptive policies must be nonanticipative. A robust criterion is an editorial choice unless attributed explicitly.

\textbf{Scope overlaps.} Allocation, collective choice, contracts, household bargaining and coordinated policy overlap economics with optimization and game theory. They are retained because they appeared in the original catalogue; no claim of a strictly game-theory-free selection is made.
% END ECONOMICS STATEMENT
\end{document}
